\documentclass[12pt]{article}

\usepackage{sbc-template}
\usepackage{graphicx,url}
\usepackage[utf8]{inputenc}
\usepackage[brazil]{babel}
\usepackage{amsmath}
\usepackage{booktabs}
\usepackage{multirow}
\usepackage{subcaption}
\usepackage{wrapfig}

\sloppy

\title{Remoção de Ruído Sal e Pimenta em Imagens de Alta Resolução: \\ Uma Análise Comparativa entre Implementações Sequencial e Paralela com OpenMP}

\author{Kauan da Rosa Paulino\inst{1}}

\address{Bacharelado em Ciência da Computação\\
Instituto Federal do Paraná (IFPR) -- Campus Pinhais\\
20251pin10020017@estudantes.ifpr.edu.br}

\begin{document}

\maketitle

\begin{resumo}
	Este trabalho investiga o desempenho de um filtro espacial da mediana $5 \times 5$ aplicado à remoção de ruído do tipo \textit{sal e pimenta} em imagens digitais de alta resolução. Foram implementadas duas versões em linguagem C --- uma puramente sequencial e outra paralelizada com OpenMP, utilizando-se o paradigma de paralelismo de dados via \texttt{\#pragma omp parallel for} com escalonamento estático. Os experimentos cobriram quatro imagens, de $0{,}26$ a $632$ megapixels, com foco nas entradas Grande ($7854 \times 11498$, $\approx 90$ Mpixels) e Extra Grande ($30000 \times 21059$, $\approx 632$ Mpixels), avaliando-se as métricas de \textit{speedup}, eficiência e \textit{overhead} para $2$, $4$, $8$, $12$ e $16$ threads, além do \textit{baseline} sequencial. Os resultados são analisados à luz das Leis de Amdahl e de Gustafson-Barsis, evidenciando comportamento próximo do linear até $4$ threads, ganho superlinear em $p=2$ atribuído a efeitos de cache e saturação a partir de $8$ threads, consistente com a topologia física do processador utilizado ($8$ núcleos, $16$ threads lógicas).
\end{resumo}

\begin{abstract}
	This work investigates the performance of a $5 \times 5$ median spatial filter applied to salt-and-pepper noise removal in high-resolution digital images. Two C implementations were developed --- a purely sequential one and one parallelized with OpenMP using data parallelism via \texttt{\#pragma omp parallel for} with static scheduling. Experiments covered four images ranging from $0.26$ to $632$ megapixels, with focus on the Large ($7854 \times 11498$, $\approx 90$ Mpixels) and Extra-Large ($30000 \times 21059$, $\approx 632$ Mpixels) inputs, evaluating \textit{speedup}, efficiency, and \textit{overhead} for $2$, $4$, $8$, $12$, and $16$ threads, in addition to the sequential baseline. Results are analyzed under Amdahl's and Gustafson-Barsis' laws, showing near-linear behavior up to $4$ threads, a superlinear gain at $p=2$ attributed to cache effects, and saturation beyond $8$ threads, consistent with the physical topology of the CPU ($8$ cores, $16$ logical threads).
\end{abstract}

\newpage
\section{Introdução}

A remoção de ruído é uma etapa fundamental em pipelines de processamento digital de imagens, com aplicações diretas em imageamento médico, sensoriamento remoto, restauração de acervos fotográficos e pré-processamento para algoritmos de visão computacional \cite{gonzalez2010processamento}. Dentre os tipos de ruído mais comuns, destaca-se o \textit{Sal e Pimenta} (\textit{impulse noise}), caracterizado pela substituição aleatória de pixels por valores extremos (0 ou 255 em imagens de 8 bits), o que degrada severamente a informação visual.

O filtro espacial da mediana é reconhecidamente uma das técnicas mais eficazes para atenuar esse tipo de ruído, pois preserva bordas melhor do que filtros lineares como a média aritmética \cite{gonzalez2010processamento}. Entretanto, sua natureza não linear — cada pixel de saída requer a ordenação de uma janela de vizinhança — o torna computacionalmente custoso, especialmente em imagens de altíssima resolução. Esse custo cresce com o quadrado da dimensão da janela ($k^2$), tornando a execução sequencial proibitiva em cenários com restrição temporal.

Nesse contexto, a paralelização via OpenMP apresenta-se como uma estratégia natural: o filtro é um exemplo clássico de \textit{paralelismo de dados embaraçosamente paralelo}, uma vez que cada pixel de saída depende exclusivamente de uma região da imagem de entrada (que não é modificada durante o processamento), conforme discutido em \cite{eijkhout2022arthpc2} e \cite{burkardt2017imagedenoise}.

Este trabalho tem como objetivo implementar, medir e analisar criticamente uma versão sequencial e uma versão paralela do filtro da mediana $5 \times 5$ aplicado a uma imagem de aproximadamente 90 megapixels, comparando o desempenho observado com as previsões teóricas das Leis de Amdahl \cite{eijkhout2022arthpc1} e de Gustafson-Barsis. Como contribuição adicional, discute-se a influência do \textit{overhead} de paralelização, dos efeitos de cache e da topologia física do processador sobre a escalabilidade do problema.

\newpage
\section{Fundamentação Teórica}

\subsection{Métricas de Desempenho em Computação Paralela}

A avaliação quantitativa de uma implementação paralela exige métricas que capturem não apenas o desempenho absoluto, mas a efetividade do uso dos recursos. Seguindo a notação de \cite{eijkhout2022arthpc1}, define-se:

\begin{itemize}
	\item \textbf{Tempo sequencial} ($T_s$): tempo de execução do melhor algoritmo sequencial disponível para resolver o problema com uma dada entrada.
	\item \textbf{Tempo paralelo} ($T_p(p)$): tempo de execução do algoritmo paralelo com $p$ unidades de processamento, para a mesma entrada.
	\item \textbf{\textit{Speedup}} ($S(p)$): $S(p) = T_s / T_p(p)$.
	\item \textbf{Eficiência} ($E(p)$): $E(p) = S(p) / p$, medindo o aproveitamento relativo dos recursos.
	\item \textbf{\textit{Overhead}} ($T_o$): $T_o(p) = T_p(p) \cdot p - T_s$, quantificando o custo adicional introduzido pela paralelização (criação/gerenciamento de threads, sincronização, comunicação, computação redundante).
\end{itemize}

Valores de $S(p) > p$ (denominados \textit{speedup superlinear}) não são previstos pela teoria clássica, mas podem manifestar-se na prática devido a efeitos de hierarquia de memória: ao particionar o \textit{working set} entre threads, cada uma passa a acessar uma região menor, potencialmente residente em caches L2/L3, reduzindo a latência média de acesso \cite{eijkhout2022arthpc1}.

\subsection{Leis de Amdahl e Gustafson-Barsis}

A Lei de Amdahl \cite{eijkhout2022arthpc1} estabelece que o \textit{speedup} de um programa com fração sequencial $\beta$ ($0 \le \beta \le 1$) é limitado por:
\begin{equation}
	S(p) = \frac{1}{\beta + \frac{1 - \beta}{p}}
\end{equation}
com limite assintótico $S_{\max} = 1/\beta$ quando $p \to \infty$. Essa lei pressupõe \textbf{tamanho de problema fixo} (escalabilidade forte).

A Lei de Gustafson-Barsis, por sua vez, assume que o tamanho do problema cresce proporcionalmente ao número de processadores (escalabilidade fraca). Seja $\alpha$ a fração sequencial no tempo paralelo:
\begin{equation}
	S(p) = \alpha + p(1 - \alpha)
\end{equation}
prevendo \textit{speedup} praticamente linear para problemas que crescem com $p$, o que reflete cenários reais de \textit{big data} e simulação científica \cite{eijkhout2022arthpc1}.

\newpage
\section{Solução Sequencial}

\subsection{Algoritmo}

O algoritmo implementado opera em imagens no formato binário \textit{Portable PixMap} (PPM P6), cada pixel representado por três canais de 8 bits (R, G, B). Para cada pixel interno $(y, x)$ da imagem, com $y \in [2, H - 3]$ e $x \in [2, W - 3]$:

\begin{enumerate}
	\item Coleta-se a janela de vizinhança $5 \times 5$ centrada em $(y, x)$, para cada canal de cor independentemente, totalizando 25 valores.
	\item Ordena-se o vetor resultante com \textit{Selection Sort}.
	\item Escreve-se o elemento mediano (índice 12, correspondente ao 13º menor valor) no pixel correspondente da imagem de saída.
\end{enumerate}

As bordas da imagem (duas linhas/colunas externas) são preservadas inalteradas por cópia direta, evitando a necessidade de \textit{padding} e mantendo a mesma dimensão da entrada.

A escolha da janela $5 \times 5$ em detrimento de $3 \times 3$ justifica-se por dois motivos: (i) maior efetividade na remoção de ruído impulsivo denso, conforme discutido em \cite{gonzalez2010processamento}; e (ii) aumento do custo computacional por pixel em $\approx 7{,}7\times$ (de 81 para 625 comparações por canal), necessário para atender ao requisito de que o programa sequencial demore ``vários minutos'' em entradas grandes.

\subsection{Decisões de Projeto}

\begin{itemize}
	\item \textbf{Selection Sort} foi escolhido por ser determinístico, sem alocação dinâmica e com \textit{overhead} de chamada de função desprezível para $n=25$. Alternativas como \textit{Quick Sort} introduziriam recursão e pior comportamento de cache para vetores pequenos.
	\item \textbf{Leitura/escrita via \texttt{fread}/\texttt{fwrite}} sobre o \textit{buffer} completo da imagem, minimizando chamadas de sistema.
	\item \textbf{Cópia integral de \texttt{img\_in} para \texttt{img\_out}} antes da filtragem, garantindo preservação das bordas com um único \texttt{memcpy} e evitando ramificações condicionais dentro do laço principal.
	\item \textbf{Medição de tempo com \texttt{omp\_get\_wtime()}} restrita ao laço de filtragem, conforme exigido pela especificação. As etapas de I/O e alocação são deliberadamente excluídas da região cronometrada, pois não fazem parte da região paralelizável.
\end{itemize}

\newpage
\section{Estratégia de Paralelização}

A paralelização baseia-se no paradigma de \textit{paralelismo de dados}: cada iteração do laço mais externo (linha $y$ da imagem) é atribuída a uma thread, conforme o trecho abaixo:

\begin{verbatim}
#pragma omp parallel for schedule(static) default(none) \
        shared(img_in, img_out, larg, alt)
for (int y = 2; y < alt - 2; y++) { ... }
\end{verbatim}

Justificam-se as escolhas de cláusulas:

\begin{itemize}
	\item \textbf{Laço externo (\texttt{y}) paralelizado}, e não o interno (\texttt{x}): cada thread escreve em uma faixa horizontal distinta da imagem de saída, eliminando a possibilidade de \textit{false sharing} — múltiplas threads escrevendo em posições de memória distintas mas adjacentes à mesma linha de cache, que causaria invalidação de cache por \textit{ping-pong} de coerência.
	\item \textbf{\texttt{schedule(static)}}: a carga por linha é uniforme (cada pixel custa o mesmo número de operações), de modo que o escalonamento estático elimina o \textit{overhead} de decisão em tempo de execução presente em \texttt{dynamic}.
	\item \textbf{\texttt{default(none)}}: obriga a explicitação do escopo de todas as variáveis, prevenindo erros sutis e documentando a intenção do autor.
	\item \textbf{\texttt{shared}} apenas para ponteiros e dimensões: as variáveis \texttt{janela}, \texttt{idx} e as de iteração são naturalmente privadas por escopo, e o compilador garante essa privacidade.
\end{itemize}

A região paralelizável exclui deliberadamente as operações de leitura, alocação e escrita de arquivo, que são inerentemente sequenciais e ficaram fora da região cronometrada.

\newpage
\section{Metodologia Experimental}

\subsection{Hardware e Software}

Os experimentos foram conduzidos em um \textit{desktop} com as seguintes características:

\begin{itemize}
	\item \textbf{CPU}: AMD Ryzen 7 5700X, 8 núcleos físicos, 16 threads lógicas (SMT habilitado), 3,4\,GHz base, \textit{boost} até 4,66\,GHz.
	\item \textbf{Memória}: 24\,GB DDR4-2666.
	\item \textbf{Sistema Operacional}: Linux Mint 22.1 (kernel padrão da distribuição).
	\item \textbf{Compilador}: GCC com suporte a OpenMP.
	\item \textbf{Flags de compilação}: \texttt{-std=c11 -Wall -Wextra -O3 -fopenmp}.
	\item \textbf{\textit{Governor} de CPU}: fixado em \texttt{performance} durante todos os experimentos, com o sistema sem carga concorrente (apenas terminal e editor de texto abertos).
\end{itemize}

\subsection{Entradas de Teste}

As imagens foram sintetizadas a partir de originais em domínio público e domínios públicos especializados, com ruído \textit{sal e pimenta} adicionado por script Python próprio (\texttt{generate\_noisy\_ppm.py}) e conversão para PPM P6:

\begin{table}[ht]
	\centering
	\caption{Conjunto de entradas de teste}
	\label{tab:entradas}
	\begin{tabular}{lccc}
		\toprule
		\textbf{Tamanho} & \textbf{Dimensões}   & \textbf{Pixels (M)} & \textbf{Uso principal}       \\
		\midrule
		Pequena          & $512 \times 512$     & 0,26                & Validação de corretude       \\
		Média            & $3648 \times 3648$   & 13,3                & Corretude + escalabilidade   \\
		Grande           & $7854 \times 11498$  & 90,3                & \textbf{Benchmark principal} \\
		Extra Grande     & $30000 \times 21059$ & 631,8               & Escalabilidade fraca         \\
		\bottomrule
	\end{tabular}
\end{table}

\subsection{Protocolo de Medição}

Para cada configuração $(p, \text{entrada})$, com $p \in \{1, 2, 4, 8, 12, 16\}$:

\begin{enumerate}
	\item Executaram-se \textbf{5 repetições} do programa, calculando-se a média aritmética do tempo reportado pelo timer interno (\texttt{omp\_get\_wtime}).
	\item Todos os testes foram realizados na mesma máquina, com o \textit{governor} fixado em \texttt{performance}, sem processos pesados concorrentes.
	\item A corretude foi verificada por comparação byte-a-byte (\texttt{cmp}) entre as saídas sequencial e paralela.
	\item O número de threads foi controlado pela variável de ambiente \texttt{OMP\_NUM\_THREADS}, garantindo o uso de um único binário para todos os testes.
\end{enumerate}

\newpage
\section{Resultados e Análise}

\subsection{Corretude}

A verificação byte-a-byte indicou identidade exata entre as saídas sequencial e paralela para as entradas Pequena, Média, Grande e Extra Grande (esta última com $p=4$), confirmando a ausência de \textit{data races} e de inconsistências numéricas.

\begin{figure}[ht]
	\centering
	\begin{subfigure}[b]{0.40\textwidth}
		\centering
		\includegraphics[width=\textwidth]{figuras/entrada_ruidosa.png}
		\caption{Imagem com ruído sal e pimenta sintetizado. Original em \cite{rajneesh231_salt_pepper_2024}.}
		\label{fig:antes}
	\end{subfigure}
	\hfill
	\begin{subfigure}[b]{0.40\textwidth}
		\centering
		\includegraphics[width=\textwidth]{figuras/saida_filtrada.png}
		\caption{Imagem após aplicação do filtro da mediana $5 \times 5$.}
		\label{fig:depois}
	\end{subfigure}
	\caption{Comparação visual antes e depois da filtragem. O filtro da mediana preserva bordas e remove satisfatoriamente o ruído impulsivo.}
	\label{fig:comparacao}
\end{figure}

O resultado de filtragem é visualmente satisfatório, preservando bordas e detalhes finos, conforme ilustrado na Figura~\ref{fig:comparacao} e em um recorte da Figura~\ref{fig:karl}. A ausência de artefatos visuais ou distorções confirma a integridade do algoritmo implementado.

\begin{figure}[ht]
	\centering
	\begin{subfigure}[b]{0.40\textwidth}
		\centering
		\includegraphics[width=\textwidth]{figuras/entrada_ruidosa_2.png}
		\caption{Imagem com ruído sal e pimenta sintetizado.}
		\label{fig:karlantes}
	\end{subfigure}
	\hfill
	\begin{subfigure}[b]{0.40\textwidth}
		\centering
		\includegraphics[width=\textwidth]{figuras/saida_filtrada_2.png}
		\caption{Imagem após aplicação do filtro da mediana $5 \times 5$.}
		\label{fig:karldepois}
	\end{subfigure}
	\caption{Comparação visual entre dois recortes da obra \textit{The Last Day of Pompeii} de Karl Brullov \cite{wikimedia_last_day_pompeii}, antes e depois da filtragem. A preservação de detalhes e a remoção do ruído são evidentes.}
	\label{fig:karl}
\end{figure}

\subsection{Tempos e \textit{Speedup}}

A Tabela~\ref{tab:tempos_grande} consolida as medições para a imagem Grande, e a Tabela~\ref{tab:tempos_extra} para a Extra Grande, que são as entradas de interesse principal deste trabalho. As Tabelas~\ref{tab:tempos_pequena} e \ref{tab:tempos_media} apresentam as entradas menores como referência complementar.

\begin{table}[ht]
	\centering
	\caption{Métricas para a imagem Grande ($7854 \times 11498$), $T_s = 86{,}606$\,s.}
	\label{tab:tempos_grande}
	\begin{tabular}{ccccc}
		\toprule
		$p$     & $T_p(p)$ (s) & $S(p)$ & $E(p)$ & $T_o(p)$ (s) \\
		\midrule
		1 (seq) & 86,606       & 1,000  & 1,000  & 0,000        \\
		2       & 42,276       & 2,049  & 1,024  & $-2,055$     \\
		4       & 21,688       & 3,993  & 0,998  & 0,146        \\
		8       & 11,909       & 7,272  & 0,909  & 8,666        \\
		12      & 10,675       & 8,113  & 0,676  & 41,496       \\
		16      & 9,052        & 9,568  & 0,598  & 58,224       \\
		\bottomrule
	\end{tabular}
\end{table}

\begin{table}[ht]
	\centering
	\caption{Métricas para a imagem Extra Grande ($30000 \times 21059$), $T_s = 616{,}853$\,s.}
	\label{tab:tempos_extra}
	\begin{tabular}{ccccc}
		\toprule
		$p$     & $T_p(p)$ (s) & $S(p)$ & $E(p)$ & $T_o(p)$ (s) \\
		\midrule
		1 (seq) & 616,853      & 1,000  & 1,000  & 0,000        \\
		2       & 299,984      & 2,056  & 1,028  & $-16,885$    \\
		4       & 156,763      & 3,935  & 0,984  & 10,199       \\
		8       & 84,851       & 7,270  & 0,909  & 61,953       \\
		12      & 76,067       & 8,109  & 0,676  & 295,946      \\
		16      & 64,291       & 9,595  & 0,600  & 411,807      \\
		\bottomrule
	\end{tabular}
\end{table}

\begin{table}[ht]
	\centering
	\caption{Médias e métricas para a imagem Pequena ($512 \times 512$), $T_s = 0{,}249$\,s.}
	\label{tab:tempos_pequena}
	\begin{tabular}{ccccc}
		\toprule
		$p$     & $T_p(p)$ (s) & $S(p)$ & $E(p)$ & $T_o(p)$ (s) \\
		\midrule
		1 (seq) & 0,249        & 1,000  & 1,000  & 0,000        \\
		2       & 0,119        & 2,090  & 1,045  & $-0,011$     \\
		4       & 0,061        & 4,093  & 1,023  & $-0,006$     \\
		8       & 0,039        & 6,464  & 0,808  & 0,059        \\
		12      & 0,032        & 7,809  & 0,651  & 0,134        \\
		16      & 0,028        & 8,831  & 0,552  & 0,202        \\
		\bottomrule
	\end{tabular}
\end{table}

\begin{table}[ht]
	\centering
	\caption{Médias e métricas para a imagem Média ($3648 \times 3648$), $T_s = 11{,}859$\,s.}
	\label{tab:tempos_media}
	\begin{tabular}{ccccc}
		\toprule
		$p$     & $T_p(p)$ (s) & $S(p)$ & $E(p)$ & $T_o(p)$ (s) \\
		\midrule
		1 (seq) & 11,859       & 1,000  & 1,000  & 0,000        \\
		2       & 5,624        & 2,109  & 1,054  & $-0,611$     \\
		4       & 2,930        & 4,048  & 1,012  & $-0,141$     \\
		8       & 1,591        & 7,455  & 0,932  & 0,866        \\
		12      & 1,475        & 8,040  & 0,670  & 5,841        \\
		16      & 1,250        & 9,489  & 0,593  & 8,138        \\
		\bottomrule
	\end{tabular}
\end{table}




\subsection{Fração Sequencial e Lei de Amdahl}

A partir da diferença sistemática entre o tempo de parede total ($T_{\text{real}}$,
medido com \texttt{/usr/bin/time}) e o tempo cronometrado da região paralelizável
($T_s$, medido com \texttt{omp\_get\_wtime}), estima-se a fração sequencial efetiva:

\begin{equation}
	\beta \approx \frac{T_{\text{real}} - T_s}{T_{\text{real}}}
\end{equation}

Para a imagem Grande, obteve-se $T_s \approx 86{,}61$\,s e
$T_{\text{real}} \approx 87{,}31$\,s, resultando em $\beta \approx 0{,}80\%$.
Para a Extra Grande, $T_s \approx 616{,}85$\,s e $T_{\text{real}} \approx 622{,}17$\,s,
fornecendo $\beta \approx 0{,}85\%$. A consistência entre as duas estimativas
indica que a fração sequencial é dominada por operações de E/S e alocação,
que crescem proporcionalmente ao tamanho da imagem.

O limite teórico de Amdahl é, portanto, $S_{\max} \approx 1/\beta \approx 120\times$,
valor inalcançável na prática: o hardware possui apenas 8 núcleos físicos, e
mesmo com 16 threads lógicas o \textit{speedup} observado satura em $\approx 9{,}6\times$.
Esse diagnóstico é fundamental: a fração sequencial \emph{não é o gargalo real
	deste problema}; o limite prático decorre da topologia física da CPU.

\subsection{Escalabilidade}

A Figura~\ref{fig:speedup} mostra o \textit{speedup} observado em função de $p$
para as quatro imagens de teste. A curva ideal (linha tracejada) é sobreposta
para referência. A Figura~\ref{fig:eficiencia} apresenta a eficiência
correspondente, e a Figura~\ref{fig:tempos} os tempos absolutos em escala
logarítmica.

\begin{figure}[ht]
	\centering
	\includegraphics[width=0.5\textwidth]{figuras/speedup.png}
	\caption{\textit{Speedup} observado vs. ideal, em função do número de threads.}
	\label{fig:speedup}
\end{figure}

\begin{figure}[ht]
	\centering
	\includegraphics[width=0.5\textwidth]{figuras/eficiencia.png}
	\caption{Eficiência $E(p)$ por número de threads. Valores acima de 1,0
		(em $p=2$ e $p=4$) indicam \textit{speedup} superlinear.}
	\label{fig:eficiencia}
\end{figure}

\begin{figure}[ht]
	\centering
	\includegraphics[width=0.5\textwidth]{figuras/tempos.png}
	\caption{Tempos absolutos de filtragem em escala logarítmica.}
	\label{fig:tempos}
\end{figure}

Três regimes distintos são observados:

\begin{enumerate}
	\item \textbf{Superlinearidade em $p = 2$} ($S \approx 2{,}05$, $E > 1$).
	      A eficiência ligeiramente acima de 1,0 e o \textit{overhead} negativo
	      ($T_o < 0$ nas Tabelas~\ref{tab:tempos_grande} e \ref{tab:tempos_extra})
	      indicam que 2 threads processam a imagem mais do que duas vezes mais rápido
	      que 1 thread sozinha. Esse fenômeno, ainda que surpreendente, é previsto
	      teoricamente: ao dividir o \textit{working set} em duas metades, cada thread
	      passa a acessar uma região menor da imagem, potencialmente com melhor
	      localidade nos caches L2 e L3 do processador.

	\item \textbf{Comportamento próximo do linear em $p \in \{2, 4\}$}: a
	      eficiência se mantém em $\approx 1{,}00$, indicando que o paralelismo
	      até 4 threads explora quase perfeitamente os recursos.

	\item \textbf{Saturação a partir de $p = 8$}: a eficiência cai para 0,91
	      ($p=8$), 0,68 ($p=12$) e 0,60 ($p=16$). Esse comportamento é esperado em
	      processadores com SMT: apenas 8 unidades de execução físicas estão
	      disponíveis, e a partir desse ponto a adição de threads traz ganhos
	      marginais decrescentes. A CPU observada durante os testes com $p=16$
	      reportou ~1460\% de uso, ou seja, apenas 14,6 dos 16 threads lógicos
	      efetivamente trabalhando em paralelo.
\end{enumerate}

É notável que as curvas de \textit{speedup} das imagens Grande e Extra Grande
sejam praticamente idênticas, apesar de a Extra Grande ser $\approx 7\times$
maior em número de pixels. Isso confirma que, uma vez que o problema é grande
o suficiente para amortizar os custos fixos de paralelização, o
comportamento de escalabilidade passa a ser determinado pelo hardware, e não
pelo tamanho da entrada.

\subsection{Escalabilidade Fraca (Gustafson-Barsis)}

A Figura~\ref{fig:escalabilidade_fraca} apresenta a eficiência para cada
tamanho de imagem em função do número de threads. Observa-se que a eficiência
em $p = 2$ e $p = 4$ é próxima ou superior a 1,0 para todas as imagens,
enquanto em $p \geq 8$ a queda é acentuada.

\begin{figure}[ht]
	\centering
	\includegraphics[width=0.5\textwidth]{figuras/escalabilidade_fraca.png}
	\caption{Eficiência por tamanho de imagem. Curvas mais altas em tamanhos
		maiores corroboram a previsão de Gustafson-Barsis.}
	\label{fig:escalabilidade_fraca}
\end{figure}

A previsão de Gustafson-Barsis é corroborada pelos dados: embora a
eficiência em $p=16$ seja baixa (0,60), o \textit{speedup} absoluto em
$p=16$ foi de 9,57$\times$ para a Grande e 9,60$\times$ para a Extra
Grande. Em outras palavras, aumentar o problema junto com o número de
threads permite ganhos próximos do limite físico do hardware, sem o
``teto'' rígido previsto pela Lei de Amdahl. A diferença entre as duas leis
está, portanto, não em qual é ``correta'', mas em qual premissa reflete o
cenário de uso: problema fixo vs. problema crescente.

\subsection{Discussão: a Solução é Interessante?}

Do ponto de vista prático, a resposta é claramente afirmativa:

\begin{itemize}
	\item \textbf{Paralelização trivial e efetiva.} O filtro da mediana é um
	      exemplo canônico de paralelismo de dados: cada pixel de saída depende
	      apenas de uma vizinhança da imagem de entrada. Não há dependências entre
	      iterações, o que permite paralelizá-lo com uma única diretiva OpenMP e
	      obter \textit{speedup} próximo do linear até o limite físico da CPU.

	\item \textbf{Ganho de tempo expressivo.} A imagem Extra Grande ($\approx 632$
	      megapixels) leva cerca de 10 minutos para ser filtrada sequencialmente,
	      mas apenas 64 segundos com 16 threads --- uma redução de $\approx 10\times$
	      no tempo de resposta.

	\item \textbf{Limitação identificada: saturação por SMT.} A curva de
	      eficiência deixa claro que, para este hardware, o ponto ótimo de
	      escalabilidade está em $p = 4$ (eficiência $\approx 1{,}00$). Ir além
	      disso traz ganhos decrescentes. Em ambientes com restrição energética ou
	      com muitos processos concorrentes, usar 16 threads seria desperdício.

	\item \textbf{Limitação algorítmica: complexidade da ordenação.} Cada
	      pixel requer a ordenação de 25 valores (Selection Sort, $O(k^4)$ no total).
	      Alternativas como redes de ordenação (Batcher's odd-even merge) ou
	      \textit{quickselect} poderiam reduzir o custo por pixel, ao preço de maior
	      complexidade de implementação. Para janelas muito grandes ($k \geq 9$),
	      essa otimização seria mandatória.
\end{itemize}

Em suma, a solução proposta é adequada ao domínio de processamento de
imagens em lote, onde múltiplas imagens de alta resolução precisam ser
processadas em tempo razoável, e onde a presença de múltiplos núcleos é
explorável de forma quase automática via OpenMP.

\newpage
\section{Conclusão}

Este trabalho implementou e avaliou duas versões de um filtro da mediana $5 \times 5$ para remoção de ruído \textit{sal e pimenta} em imagens de alta resolução. A versão paralela com OpenMP demonstrou \textit{speedup} substancial sobre a sequencial, com comportamento próximo do linear até o limite físico da CPU utilizada. A fração sequencial estimada ($\beta \approx 0{,}83\%$) indica que o problema é altamente paralelizável, e o limite prático de escalabilidade decorre da topologia do processador, não do algoritmo. Os resultados confirmam a adequação do paradigma de paralelismo de dados a este domínio e servem como caso de estudo da aplicabilidade das Leis de Amdahl e Gustafson-Barsis.

Os resultados confirmam a adequação do paradigma de paralelismo de dados a este domínio. A eficiência próxima da unidade até $p=4$ e a saturação a partir de $p=8$ evidenciam que o limite prático de escalabilidade é imposto pela topologia física do processador (8 núcleos), e não pelo algoritmo. A fração sequencial estimada ($\beta \approx 0{,}8\%$) coloca o limite de Amdahl em $\approx 120\times$, valor irrealizável neste hardware e que ilustra bem a diferença entre o limite teórico e o limite prático da escalabilidade forte.

\newpage
\bibliographystyle{sbc}
\bibliography{referencias}

\end{document}